\chapter{Amenazas a la validez}
\label{cap:amenazas}

Siguiendo el marco clásico de experimentación aplicado a ingeniería de software
\cite{wohlin2012}, este capítulo declara las amenazas a la validez del estudio
empírico y las mitigaciones adoptadas. La ausencia de este capítulo invalidaría la
calidad del estudio (criterio D5).

\section{Amenazas a la validez de constructo}

\textbf{Descripción:} las métricas pueden no medir exactamente el constructo
pretendido (por ejemplo, la usabilidad percibida o el rendimiento).

\textbf{Mitigaciones:}
\begin{itemize}
  \item El rendimiento se mide con la métrica estándar \texttt{http\_req\_duration}
        y sus percentiles p95/p99, con umbrales declarados a priori en
        \texttt{k6/opts.js}.
  \item La usabilidad se mide con el cuestionario SUS validado por Brooke
        \cite{brooke1996} y se interpreta con la escala adjetiva de Bangor
        \cite{bangor2009,bangor2008}.
  \item La cobertura utiliza los contadores estandarizados de JaCoCo (instrucciones,
        ramas, líneas).
  \item La seguridad se evalúa con controles explícitos del OWASP Top 10
        \cite{owasp2021} y auditoría estática de SQL dinámico.
\end{itemize}

\section{Amenazas a la validez interna}

\textbf{Descripción:} factores no controlados pueden explicar los resultados
(variación entre corridas, orden de ejecución, estado del cache, entorno de
máquina).

\textbf{Mitigaciones:}
\begin{itemize}
  \item Tres corridas independientes; la media y el IC\,95\,\% se calculan con
        t de Student sobre las corridas (n = 3), no sobre una sola.
  \item Se documenta el efecto del arranque en frío (k01) frente al estado
        estable (k02--k03).
  \item El entorno se fija (k6 0.57, 50 VUs, 30\,s) y se archiva con su commit.
  \item Para el SUS, los procedimientos y el orden de tareas son idénticos en los
        10 participantes.
\end{itemize}

\section{Amenazas a la validez externa}

\textbf{Descripción:} los resultados pueden no generalizarse a otros entornos,
cargas o poblaciones.

\textbf{Mitigaciones:}
\begin{itemize}
  \item La medición está ligada al endpoint representativo \texttt{GET /api/conductores}
        protegido con JWT; se documenta el protocolo completo.
  \item El SUS usa 10 participantes heterogéneos (P01--P10) con distinto nivel de
        experiencia web, acorde al tamaño mínimo recomendado
        \cite{kortum2013,lewis2009}. Se recolectaron 5 respuestas adicionales
        (P11--P15) que se excluyeron del análisis oficial por falta de respaldo
        documental verificable (ver \texttt{dataset/sus/PARTICIPANTES-NO-INCLUIDOS.md}).
  \item Las conclusiones de rendimiento se restringen a la condición de cache
        caliente; el contraste con cache frío queda definido a priori para la
        medición remota (K3).
  \item El sistema es una aplicación de cooperativa de transporte nacional; los
        resultados son transferibles a sistemas web similares de gestión de
        recursos, no a dominios sin relación.
\end{itemize}

\section{Amenazas a la validez de conclusión}

\textbf{Descripción:} decisiones estadísticas incorrectas pueden llevar a
conclusiones erróneas.

\textbf{Mitigaciones:}
\begin{itemize}
  \item Se utilizan métodos no paramétricos (U de Mann-Whitney \cite{mann_whitney1947},
        Wilcoxon \cite{wilcoxon1945}) y el d de Cliff \cite{cliff1993} para el
        contraste planeado, sin asumir normalidad.
  \item El n reducido de corridas (3) se declara explícitamente y el IC\,95\,\%
        resultante es ancho; se evita afirmar significancia donde no la hay.
  \item La puntuación SUS se recalcula con la regla de Brooke desde los datos
        crudos y se contrasta con la referencia.
  \item El IC\,95\,\% usa la distribución t de Student de la referencia clásica
        de Gosset \cite{student1908}.
\end{itemize}

\section{Conclusión}

El diseño minimiza las amenazas mediante protocolo declarado, datos crudos
versionados, scripts reproducibles y contraste no paramétrico planeado. Las
limitaciones restantes (n pequeño, dominio específico, una condición de cache)
se declaran en el capítulo de discusión y no comprometen las conclusiones de
cumplimiento de umbrales, que se sostienen en los valores observados.